
\documentclass[11pt]{article}

\usepackage[letterpaper,margin=0.78in]{geometry}
\usepackage[T1]{fontenc}
\usepackage[utf8]{inputenc}
\usepackage{lmodern}
\usepackage{microtype}
\usepackage[table]{xcolor}
\usepackage{booktabs}
\usepackage{tabularx}
\usepackage{longtable}
\usepackage{array}
\usepackage{enumitem}
\usepackage{hyperref}
\usepackage{xurl}
\usepackage{fancyhdr}
\usepackage{titlesec}
\usepackage[most]{tcolorbox}

\definecolor{accent}{HTML}{1F4E79}
\definecolor{lightaccent}{HTML}{EAF2F8}
\definecolor{midgray}{HTML}{666666}
\definecolor{lightgray}{HTML}{F5F7FA}
\definecolor{warn}{HTML}{8A4B08}
\definecolor{success}{HTML}{2E7D32}
\definecolor{softgreen}{HTML}{EAF7EA}

\hypersetup{
  colorlinks=true,
  linkcolor=accent,
  urlcolor=accent,
  citecolor=accent,
  pdftitle={JAN-PRO Franchise Sales Director Meeting Prep for Jordan Suber},
  pdfauthor={Jordan Suber},
  pdfsubject={Tailored JAN-PRO franchise candidate meeting preparation},
  pdfkeywords={JAN-PRO, Jordan Suber, franchise, commercial cleaning, Franchise Disclosure Document, FDD, SBA, FTC, due diligence}
}

\setlength{\headheight}{14pt}
\pagestyle{fancy}
\fancyhf{}
\lhead{JAN-PRO Franchise Meeting Prep}
\rhead{Jordan Suber Candidate Handout}
\cfoot{\thepage}
\renewcommand{\headrulewidth}{0.4pt}

\titleformat{\section}{\Large\bfseries\color{accent}}{}{0em}{}[\titlerule]
\titleformat{\subsection}{\large\bfseries\color{accent}}{}{0em}{}
\titleformat{\subsubsection}{\normalsize\bfseries\color{accent}}{}{0em}{}
\setlist[itemize]{leftmargin=1.35em,itemsep=0.22em,topsep=0.25em}
\setlist[enumerate]{leftmargin=1.45em,itemsep=0.22em,topsep=0.25em}
\renewcommand{\arraystretch}{1.22}
\emergencystretch=2em
\sloppy

\newtcolorbox{keybox}[1]{
  colback=lightaccent,
  colframe=accent,
  title=\textbf{#1},
  arc=2mm,
  boxrule=0.6pt,
  left=1.1mm,
  right=1.1mm,
  top=0.9mm,
  bottom=0.9mm
}

\newtcolorbox{scriptbox}[1]{
  colback=lightgray,
  colframe=midgray,
  title=\textbf{#1},
  arc=1.5mm,
  boxrule=0.45pt,
  left=1.1mm,
  right=1.1mm,
  top=0.9mm,
  bottom=0.9mm
}

\newtcolorbox{warningbox}[1]{
  colback=orange!7,
  colframe=warn,
  title=\textbf{#1},
  arc=1.5mm,
  boxrule=0.45pt,
  left=1.1mm,
  right=1.1mm,
  top=0.9mm,
  bottom=0.9mm
}

\newtcolorbox{greenbox}[1]{
  colback=softgreen,
  colframe=success,
  title=\textbf{#1},
  arc=1.5mm,
  boxrule=0.45pt,
  left=1.1mm,
  right=1.1mm,
  top=0.9mm,
  bottom=0.9mm
}

\newcommand{\placeholder}[1]{\textcolor{accent}{\texttt{[#1]}}}
\newcommand{\janpro}{JAN-PRO}
\newcommand{\candidate}{Jordan Suber}

\begin{document}

\begin{center}
  {\LARGE\bfseries Meeting with a \janpro{} Director of Franchise Sales}\\[0.35em]
  {\large Tailored Candidate Positioning Handout for \candidate}\\[0.4em]
  {\small Commercial-cleaning franchise discovery, qualification, and due-diligence preparation}
\end{center}

\vspace{0.5em}

\begin{keybox}{Jordan's meeting objective}
This meeting should position \candidate{} as a serious franchise candidate with a technology, DevSecOps, cloud architecture, operations, automation, governance, and audit-readiness background. The goal is not to sound like a passive investor or a person simply buying a cleaning route. The goal is to communicate that Jordan can evaluate and operate a recurring B2B service business using disciplined process management, customer follow-up, measurable quality control, staff accountability, conservative financial diligence, and a willingness to follow a franchise system.
\end{keybox}

\begin{warningbox}{Important financial boundary}
The only facts that should be stated as final are facts Jordan can verify. Use placeholders for liquid capital, net worth, financing plan, and investment ceiling until those figures are confirmed. Do not improvise numbers during the meeting. A franchise sales director will expect directness, but directness is different from overstatement.
\end{warningbox}

\section{Candidate Snapshot: Jordan Suber}

\begin{longtable}{>{\raggedright\arraybackslash}p{0.31\linewidth}>{\raggedright\arraybackslash}p{0.63\linewidth}}
\toprule
\textbf{Candidate dimension} & \textbf{How to frame it in the meeting} \\
\midrule
Professional identity & Software engineer / senior enterprise cloud architect with strong DevSecOps, infrastructure, platform engineering, automation, security controls, process governance, and operational review experience. \\
Operating strengths & Builds repeatable systems, documents procedures, manages risk, designs controls, automates manual workflows, and thinks in terms of service quality, evidence, metrics, accountability, and continuous improvement. \\
Business posture & Founder/operator mindset: willing to validate the opportunity, follow the system, learn the commercial-cleaning operating model, and scale only after the first operating base is stable. \\
Market posture & Atlanta Metropolitan Area / North Georgia commercial corridors are the likely starting market unless the franchisor maps Jordan to a different available territory. Replace this with the exact target market before the call. \\
Franchise fit thesis & Commercial cleaning can fit Jordan's strengths because it rewards disciplined B2B account management, service-level execution, inspection routines, workforce reliability, customer-retention discipline, and measurable quality control. \\
Primary risk to manage & The director may wonder whether a technical professional understands that this is a local service business, not a software project. Jordan should answer this directly by emphasizing owner involvement, sales discipline, field execution, and willingness to learn the system. \\
\bottomrule
\end{longtable}

\section{Jordan's Core Positioning Message}

\begin{scriptbox}{One-minute opening script}
``Thank you for meeting with me. I am Jordan Suber. My background is in software engineering, cloud architecture, DevSecOps, automation, security controls, and operational governance. I am evaluating franchise opportunities where disciplined execution, repeatable systems, measurable service quality, and customer retention matter. \janpro{} interests me because commercial cleaning is a recurring B2B service model that appears to reward process discipline, sales follow-up, staffing reliability, quality control, and consistent account management. I am here to understand whether I am a strong candidate for your system, and whether the system is the right fit for my capital, target market, timeline, and long-term ownership goals.''
\end{scriptbox}

\begin{scriptbox}{Shorter version if the call moves quickly}
``I am Jordan Suber, a technology and operations professional with a strong background in automation, security, process governance, and service reliability. I am looking for a franchise model where those strengths translate into repeatable local execution. I am interested in \janpro{} because commercial cleaning is a recurring B2B service business where account retention, quality standards, staffing reliability, and disciplined operations are critical.''
\end{scriptbox}

\begin{greenbox}{Candidate posture to maintain}
Jordan should sound like a disciplined operator-investor: direct about qualification, humble about the cleaning industry, serious about local sales and service execution, and conservative about financial claims until the FDD and franchisee validation support them.
\end{greenbox}

\section{How Jordan Should Explain Why \janpro{} Fits}

Avoid a generic answer such as ``I like the brand.'' Tie the fit to operating logic and Jordan's actual strengths.

\begin{scriptbox}{Brand-fit answer for Jordan}
``What attracts me to \janpro{} is the potential to build a recurring commercial-service business around disciplined execution. My professional background has trained me to build repeatable processes, manage operational risk, measure service quality, document procedures, and improve systems over time. I understand this is not a passive-income model. I would need to learn the cleaning operations, manage account expectations, control costs, handle staffing reliability, and retain customers. That operating challenge is why I want to evaluate the system seriously.''
\end{scriptbox}

\begin{scriptbox}{If asked why a technical background helps}
``My technical background matters because I am used to operating mission-critical systems where reliability, documentation, incident response, quality controls, and customer trust matter. I do not assume that technology experience automatically makes me qualified for commercial cleaning. But it does give me a disciplined operating style: define the standard, measure the work, inspect outcomes, correct defects, document the process, and improve repeatability.''
\end{scriptbox}

\begin{scriptbox}{If asked whether Jordan wants to be hands-on}
``I expect to be hands-on, especially early. I want to understand how accounts are sold, scoped, staffed, serviced, inspected, retained, and corrected when issues occur. My long-term goal may include building a scalable operation, but I would not try to scale until I understand the system and can run the first accounts reliably.''
\end{scriptbox}

\section{What the Director Needs to Learn Quickly}

A franchise sales director will likely evaluate whether Jordan is serious, qualified, coachable, financially realistic, sales-capable, and operationally credible. Jordan should prepare crisp answers in the following areas.

\begin{longtable}{>{\raggedright\arraybackslash}p{0.28\linewidth}>{\raggedright\arraybackslash}p{0.66\linewidth}}
\toprule
\textbf{Topic} & \textbf{Jordan's prepared answer} \\
\midrule
Candidate summary & ``I am a software engineering and cloud/DevSecOps professional with deep experience in automation, process controls, service reliability, governance, and operational execution.'' \\
Ownership model & ``I am evaluating \placeholder{owner-operator / executive operator / growth-oriented operator} models and want to understand which model is most realistic for \janpro{} in this market.'' \\
Liquid capital & ``My confirmed liquid-capital range is \placeholder{amount}. I want to confirm whether that aligns with your current qualification requirements and the full initial investment range.'' \\
Net worth and financing & ``My estimated net-worth range is \placeholder{range}. My financing plan is \placeholder{cash / SBA-backed loan / partner capital / combination / other}.'' \\
Target market & ``My likely target market is \placeholder{Atlanta metro / Cobb County / specific territory}. I want to understand available territories, account density, and what commercial segments perform well locally.'' \\
Sales readiness & ``I understand that this requires local B2B prospecting, networking, follow-up, proposals, account onboarding, and retention. I want to understand the franchisor's sales process and the owner's expected role.'' \\
Operations readiness & ``My strength is building repeatable operating systems: checklists, metrics, quality inspections, escalation paths, documentation, and corrective actions. I want to learn how \janpro{} expects those activities to be performed.'' \\
Due-diligence posture & ``I plan to review the FDD, speak with current and former franchisees, validate account economics, review the local market, and use legal/accounting review before signing.'' \\
Timeline & ``My expected decision timeline is \placeholder{exploring / 3--6 months / 6--12 months / ready to move through qualification}.'' \\
\bottomrule
\end{longtable}

\section{Regulatory and Diligence Baseline}

\begin{itemize}
  \item The Federal Trade Commission's Franchise Rule requires franchisors to provide a disclosure document containing 23 specific items about the franchise offer, officers, and other franchisees.\cite{ftc-rule}
  \item The FTC states that prospective franchise buyers must receive the FDD at least 14 days before being asked to sign a contract or pay money to the franchisor or its affiliate.\cite{ftc-consumer-guide,ftc-deep-dive}
  \item The FTC advises prospective franchise buyers to read the FDD carefully, ask questions, evaluate costs, review training and support, speak with franchisees, and treat financial-performance claims carefully.\cite{ftc-consumer-guide,ftc-deep-dive}
  \item SBA-backed financing requires brand eligibility review. SBA states that the SBA Franchise Directory contains franchises and other brands eligible for SBA financial assistance and that placement is not an endorsement or guarantee of success.\cite{sba-directory}
  \item Independent validation matters. The International Franchise Association recommends analyzing the FDD, speaking with current and former franchisees, evaluating initial and ongoing costs, researching market demand, and using independent professional advisors.\cite{ifa-diligence}
\end{itemize}

\section{Financial Qualification Statement for Jordan}

The financial statement should be confident but bounded. Jordan should be ready to discuss liquid capital, net worth, financing strategy, working-capital cushion, and a maximum all-in investment range. Do not let the conversation stay vague.

\begin{scriptbox}{Use this when asked about capital}
``I understand that franchise qualification depends on liquid capital, net worth, financing readiness, and the full initial investment range. My confirmed liquid-capital range is \placeholder{amount}, my approximate net-worth range is \placeholder{range}, and my intended funding strategy is \placeholder{cash / SBA-backed loan / partner capital / combination}. I also want to understand the working-capital cushion needed after startup costs, because I do not want to undercapitalize a recurring-service business during the account-ramp period.''
\end{scriptbox}

\begin{scriptbox}{Use this if the numbers are not final yet}
``I am still finalizing my investment ceiling, so I do not want to overstate my capital position. What I can say is that I am approaching this as a serious diligence process. Before moving beyond discovery, I want to confirm the current investment range, liquid-capital requirement, working-capital expectations, financing options, and whether the model fits my conservative operating assumptions.''
\end{scriptbox}

\subsection{Numbers Jordan should fill in before the call}

\begin{longtable}{>{\raggedright\arraybackslash}p{0.39\linewidth}>{\raggedright\arraybackslash}p{0.55\linewidth}}
\toprule
\textbf{Metric} & \textbf{Jordan's answer} \\
\midrule
Liquid capital available & \placeholder{cash, marketable securities, immediately available funds} \\
Estimated net worth & \placeholder{range} \\
Financing plan & \placeholder{cash / SBA / HELOC / investor / retirement rollover / other} \\
Target investment range & \placeholder{maximum all-in startup investment} \\
Working-capital cushion & \placeholder{months of operating cushion beyond initial investment} \\
Target market & \placeholder{Atlanta metro / Cobb County / specific city or county} \\
Ownership model & \placeholder{owner-operator, executive operator, multi-account operator} \\
Local customer focus & \placeholder{offices, medical/professional, retail, churches, schools, light industrial, other commercial facilities} \\
Decision timeline & \placeholder{exploring, 3--6 months, 6--12 months, ready now} \\
\bottomrule
\end{longtable}

\section{Likely Concerns About Jordan and Strong Answers}

\begin{longtable}{>{\raggedright\arraybackslash}p{0.34\linewidth}>{\raggedright\arraybackslash}p{0.60\linewidth}}
\toprule
\textbf{Possible concern from the Director} & \textbf{Jordan's answer} \\
\midrule
``You have a technical background. Why commercial cleaning?'' & ``I am interested because this is an operational, recurring-service business. My technical career has been about reliability, process control, documentation, governance, and service quality. I want to apply those habits to a local B2B franchise model.'' \\
``Do you understand this is not passive?'' & ``Yes. I expect to be involved in sales, account onboarding, staffing decisions, quality inspections, customer follow-up, and issue resolution, especially early.'' \\
``Can you sell?'' & ``I understand that B2B sales and follow-up are core to the model. I want to understand the expected owner role, available sales support, proposal process, local marketing expectations, and realistic sales-cycle assumptions.'' \\
``Will you follow the system?'' & ``Yes. I am not looking to reinvent the brand. I want to understand the system, follow documented standards, measure execution, and improve within the operating model.'' \\
``Are you financially ready?'' & ``I am prepared to discuss verified capital and financing assumptions. I also want to validate the full investment range and working-capital need before making any commitment.'' \\
``Are you only interested in technology?'' & ``No. Technology is one tool. The business has to work through people, service quality, accountability, customer trust, and local execution. I understand that those are the core operating realities.'' \\
\bottomrule
\end{longtable}

\section{Strategic Questions Jordan Should Ask \janpro{}}

Ask questions that sound like an operator and investor. The aim is to test the franchise structure, support model, account-acquisition process, unit economics, territory rules, labor assumptions, and candidate fit.

\subsection{Franchise structure and local support}

\begin{itemize}
  \item ``Who would be my direct franchisor or support organization in this market: the national franchisor, a regional developer, or another local franchise entity?''
  \item ``Which party provides training, account support, billing support, field inspections, complaint escalation, technology, and ongoing coaching?''
  \item ``How is the relationship between national brand standards and local/regional support managed?''
  \item ``What should I understand about the difference between a unit franchise, regional franchise, master franchise, or development opportunity if those structures apply in this market?''
  \item ``What are the most common misunderstandings candidates have about how the \janpro{} system works?''
\end{itemize}

\subsection{Account acquisition, leads, and sales expectations}

\begin{itemize}
  \item ``How are commercial accounts generated: franchisor-provided leads, regional sales support, owner prospecting, referrals, local marketing, or a combination?''
  \item ``Does the system provide initial accounts, guaranteed account volume, account packages, or lead targets? If so, where is that documented in the FDD or agreement?''
  \item ``What level of local sales activity do successful franchisees perform themselves?''
  \item ``What are the sales-cycle realities for commercial cleaning accounts in this market?''
  \item ``How does \janpro{} help franchisees prepare proposals, price work, define scope, and avoid underbidding?''
  \item ``What happens if a customer cancels, reduces scope, delays payment, or disputes service quality?''
\end{itemize}

\subsection{Territory and customer concentration}

\begin{itemize}
  \item ``What exactly is protected: geography, customer accounts, account type, lead source, or something else?''
  \item ``Can another \janpro{} franchisee serve accounts in or near my market? Under what circumstances?''
  \item ``How are national, regional, and local commercial accounts assigned?''
  \item ``What customer-concentration risk should I plan for if early revenue depends on a small number of accounts?''
  \item ``Which FDD item and agreement sections define territory, account ownership, non-compete, renewal, transfer, and termination rights?''
\end{itemize}

\subsection{Operations, quality control, and Jordan's process-management lens}

\begin{itemize}
  \item ``What does onboarding look like from signing through first account launch?''
  \item ``How much training is technical cleaning training versus sales, customer management, labor management, finance, and quality control?''
  \item ``How are cleaning standards, inspections, customer complaints, and corrective actions documented?''
  \item ``What metrics should a new franchisee watch weekly: account count, labor hours, gross margin, rework, complaints, churn, supply cost, or inspection scores?''
  \item ``What service failures most commonly hurt new franchisees, and how does \janpro{} coach owners through them?''
  \item ``What software, checklists, reporting tools, or inspection methods are required?''
  \item ``What are the expectations for night, weekend, holiday, backup, and emergency coverage?''
\end{itemize}

\subsection{Labor, compliance, insurance, and safety}

\begin{itemize}
  \item ``What labor model do successful franchisees use: owner-performed work, employees, subcontractors where permitted, or a mix?''
  \item ``What insurance, bonding, background-check, workers' compensation, licensing, or local business requirements are typical in this market?''
  \item ``How are chemical handling, safety data sheets, equipment training, OSHA expectations, and customer-site safety rules handled?''
  \item ``Are there special requirements for medical, educational, government, industrial, or regulated facilities?''
  \item ``What costs do new franchisees commonly underestimate for labor, payroll taxes, insurance, supplies, equipment replacement, vehicle use, and supervision?''
\end{itemize}

\subsection{FDD, economics, and financing}

\begin{itemize}
  \item ``Does the current FDD include Item 19 financial performance representations, and how should I interpret them conservatively?''
  \item ``Which revenue, gross margin, account-retention, cancellation, and ramp-up assumptions should I validate through franchisee calls?''
  \item ``What are the ongoing fees: royalties, brand/marketing fund, account administration, technology, training, renewal, transfer, insurance, or other required payments?''
  \item ``What is the minimum working-capital cushion you recommend beyond the stated initial investment range?''
  \item ``Do franchisees commonly use SBA-backed financing, and is the brand currently eligible in the SBA Franchise Directory?''
  \item ``Can you introduce me to lenders who understand the commercial-cleaning franchise model, without pressuring me toward one financing source?''
\end{itemize}

\section{Jordan's FDD Review Focus}

\begin{longtable}{>{\raggedright\arraybackslash}p{0.22\linewidth}>{\raggedright\arraybackslash}p{0.72\linewidth}}
\toprule
\textbf{FDD item} & \textbf{Why it matters for Jordan's diligence} \\
\midrule
Items 1--2 & Understand the franchisor, affiliates, regional structure, executives, and operating history. Confirm who actually supports Jordan locally. \\
Items 3--4 & Review litigation and bankruptcy history for franchisor, affiliates, and key executives. Look for patterns involving franchisee relationships, fees, accounts, termination, or support. \\
Items 5--7 & Build the full startup and first-year cash model: franchise fee, equipment, supplies, insurance, software, training, vehicle, payroll, local marketing, professional fees, and working capital. \\
Item 8 & Understand required suppliers, cleaning chemicals, equipment, uniforms, software, insurance, and whether the franchisor receives rebates or financial benefits from designated vendors. \\
Item 11 & Validate training, operations manuals, account support, marketing, sales assistance, technology, inspections, and field support. Ask recent franchisees whether the support matched expectations. \\
Item 12 & Read territory, account, lead, customer, and online-marketing restrictions carefully. Territory language is critical in a commercial-services model. \\
Item 17 & Review renewal, termination, transfer, dispute resolution, non-compete, and post-termination obligations before signing. \\
Item 19 & Treat financial performance representations conservatively. Ask what population of franchisees the numbers represent and whether the data reflects Jordan's target market and operating model. \\
Item 20 & Use the contact lists for current and former franchisee validation calls. Ask about account acquisition, ramp-up, support quality, cancellations, costs, and whether they would buy again. \\
Item 21 & Have an accountant review financial statements to assess whether the franchisor can deliver promised support. \\
\bottomrule
\end{longtable}

\section{Franchisee Validation Call Sheet}

Speak with current and former franchisees before making a final decision. The FTC and IFA both emphasize franchisee conversations as a key source of practical diligence.\cite{ftc-consumer-guide,ftc-deep-dive,ifa-diligence}

\begin{enumerate}
  \item How long did it take from signing to launching your first accounts?
  \item What did your actual startup cost and working-capital need look like versus the FDD range?
  \item How were your initial accounts generated, and what did you personally have to do to grow revenue?
  \item What surprised you most about labor, scheduling, inspections, supplies, insurance, or customer complaints?
  \item How responsive was local/regional support during the first 90, 180, and 365 days?
  \item What accounts were hardest to service profitably?
  \item What fees or costs did you underestimate?
  \item What is your customer-retention experience, and why do customers cancel?
  \item How many hours per week did you work in year one compared with what you expected?
  \item Knowing what you know now, would you buy this franchise again? Why or why not?
\end{enumerate}

\section{What Jordan Should Not Say}

\begin{longtable}{>{\raggedright\arraybackslash}p{0.38\linewidth}>{\raggedright\arraybackslash}p{0.56\linewidth}}
\toprule
\textbf{Avoid saying} & \textbf{Say this instead} \\
\midrule
``How much money can I make?'' & ``What financial performance representations are disclosed in Item 19, and what assumptions should I use when building a conservative pro forma?'' \\
``Is this passive income?'' & ``What level of owner involvement do your strongest franchisees typically have during the first year?'' \\
``I am a technology person, so I can automate this.'' & ``My technology background helps me think in terms of systems, metrics, documentation, and controls, but I understand the business depends first on local sales, staff reliability, and customer service.'' \\
``Will you give me accounts?'' & ``How are accounts generated, assigned, priced, onboarded, retained, and documented in the FDD or franchise agreement?'' \\
``Cleaning is simple, so this should be easy.'' & ``I understand that commercial cleaning requires sales discipline, labor reliability, quality control, customer retention, and cost management.'' \\
``I do not know my budget yet.'' & ``I am finalizing my investment ceiling and want to confirm the liquid-capital requirement, all-in startup range, and working-capital expectations before moving forward.'' \\
``I just want the cheapest route to ownership.'' & ``I want to understand the investment level that gives a new franchisee a realistic operating runway and adequate working capital.'' \\
\bottomrule
\end{longtable}

\section{Closing the Meeting}

End by forcing clarity on fit, qualification, and next steps.

\begin{scriptbox}{Professional close for Jordan}
``This has been helpful. Based on what we discussed, I would like to understand the next stage in your process. What materials should I review, what qualification information do you need from me, and what is the expected timeline from this initial conversation to application, FDD review, validation calls, discovery, and a possible franchise agreement?''
\end{scriptbox}

\begin{scriptbox}{Direct fit question}
``Based on what I have shared about my background, capital planning, market interest, and operating goals, do you see me as a strong candidate for \janpro{}? Are there any concerns about my financial profile, owner involvement, sales readiness, or operating fit that I should address before moving forward?''
\end{scriptbox}

\section{Post-Meeting Follow-Up Email Template}

\begin{scriptbox}{Follow-up email from Jordan}
\textbf{Subject:} Follow-up on \janpro{} franchise discussion\\[0.4em]
\textbf{Body:}\\
Hello \placeholder{Name},\\[0.25em]
Thank you for taking the time to speak with me about the \janpro{} franchise opportunity. I appreciated the overview of the commercial-cleaning model, candidate profile, training process, account-development expectations, local support structure, and market opportunity.\\[0.25em]
Based on our conversation, I remain interested in evaluating the opportunity further. My next diligence priorities are reviewing the current FDD, understanding the full investment range and working-capital expectations, validating the account-acquisition and support model, reviewing territory/account rules, and speaking with current and former franchisees where appropriate.\\[0.25em]
For context, my professional background is in software engineering, cloud architecture, DevSecOps, automation, operational governance, and service reliability. I am especially interested in understanding how \janpro{} supports disciplined local execution, account retention, quality control, and scalable operations.\\[0.25em]
Please send the next-step materials and any qualification documents you need from me. I am also happy to provide additional detail on my liquid capital, financing plan, target market, ownership timeline, and operating goals.\\[0.25em]
Best,\\
Jordan Suber
\end{scriptbox}

\section{Pre-Call Checklist for Jordan}

\begin{itemize}
  \item Confirm liquid-capital range and maximum all-in investment range.
  \item Confirm estimated net worth and financing plan.
  \item Decide whether the target role is owner-operator, executive operator, or future multi-account operator.
  \item Identify the exact target market: \placeholder{Atlanta metro / Cobb County / specific territory / other}.
  \item Prepare a local market thesis: commercial density, office/medical/professional corridors, local competition, commute radius, labor availability, and target account types.
  \item Prepare a concise explanation of how cloud, DevSecOps, SRE, automation, governance, and controls experience transfers to a local service business.
  \item Prepare questions that cannot be answered by the website.
  \item Prepare a question set around account generation, assignment, retention, cancellation, pricing, and support.
  \item Confirm whether SBA-backed financing is part of the plan and check the SBA Franchise Directory before relying on SBA eligibility.
  \item Prepare a post-meeting diligence plan: FDD review, franchisee calls, professional review, financing review, market review, and conservative pro forma.
\end{itemize}

\section{Red Flags to Watch For}

\begin{warningbox}{Warning signs during the sales conversation}
\begin{itemize}
  \item Pressure to pay quickly before reviewing the FDD and agreements.
  \item Informal earnings or account-volume claims that are not tied to Item 19 or written materials.
  \item Vague answers about how accounts are generated, assigned, priced, serviced, retained, or replaced after cancellation.
  \item Overemphasis on ``recurring revenue'' without clear discussion of labor, cancellations, service failures, and customer concentration.
  \item Resistance to current and former franchisee validation calls.
  \item Unclear distinction between national franchisor support and local/regional support.
  \item No direct answer on territory, account ownership, competition among franchisees, or national-account assignment.
  \item Dismissive treatment of insurance, labor law, safety, chemical handling, or site-specific compliance requirements.
\end{itemize}
\end{warningbox}

\section{Best Overall Script for Jordan}

\begin{greenbox}{Use this when a polished summary is needed}
``I am looking for a franchise opportunity where my capital, professional background, and ownership goals align with a proven operating system. My background is in software engineering, cloud architecture, DevSecOps, automation, controls, and operational governance, so I naturally look for businesses where repeatable execution, service quality, documentation, accountability, and risk management matter. I am interested in \janpro{} because commercial cleaning appears to reward disciplined B2B sales, account retention, quality control, staffing reliability, and cost management. My goal today is to understand whether I am a strong candidate for your system and whether the system is the right fit for me after FDD review, franchisee validation, unit-economics analysis, and local market diligence.''
\end{greenbox}

\renewcommand{\refname}{\color{accent}References}
\begin{thebibliography}{9}

\bibitem{ftc-rule}
Federal Trade Commission.
\textit{Franchise Rule}.
\url{https://www.ftc.gov/legal-library/browse/rules/franchise-rule}

\bibitem{ftc-consumer-guide}
Federal Trade Commission.
\textit{A Consumer's Guide to Buying a Franchise}.
\url{https://www.ftc.gov/business-guidance/resources/consumers-guide-buying-franchise}

\bibitem{ftc-deep-dive}
Federal Trade Commission.
\textit{Franchise Fundamentals: Taking a Deep Dive into the Franchise Disclosure Document}.
\url{https://www.ftc.gov/business-guidance/blog/2023/05/franchise-fundamentals-taking-deep-dive-franchise-disclosure-document}

\bibitem{sba-directory}
U.S. Small Business Administration.
\textit{SBA Franchise Directory}.
\url{https://www.sba.gov/business-guide/plan-your-business/buy-existing-business-or-franchise/sba-franchise-directory}

\bibitem{ifa-diligence}
International Franchise Association.
\textit{Diligence in Franchising}.
\url{https://www.franchise.org/franchising-overview/diligence-in-franchising/}

\end{thebibliography}

\end{document}
